\documentclass[letterpaper,twocolumn,10pt]{article}
\usepackage{hyperref}
\usepackage{usenix2019_v3_xetex}
\usepackage[T1]{fontenc}
\usepackage{amsmath,amssymb,graphicx}
\hypersetup{hidelinks,pdftitle={Mechanism figure: typesetting proof},pdfauthor={}}
\makeatletter
\newcommand{\captiontypefigure}{\def\@captype{figure}}
\makeatother
\begin{document}
\twocolumn[{
\noindent\textbf{Mechanism overview: author-review proof}\par
\vspace{8pt}
\centering
\includegraphics[width=\textwidth]{figures/system-architecture-mechanism.pdf}
\vspace{3pt}
\captiontypefigure
\caption{Private views connected to object-bound storage layouts.
Each agent resolves its pinned snapshot $S$ and private current state
$\Delta_i$. The aligned range maps illustrate different private overwrites:
$S$ cells are explicit lower references, and $U_i$ cells reference separate
private payload objects. $\Delta_i$ resides in the mutable metadata service;
$S$ and $U_i$ reside in object storage. A representative packed record
co-locates a name, common attributes and a locator to a frame in $O_1$.
The grey connector shows that same object's storage representation.
At finalization, a configured policy may request $B$; RustFS admits the layout
and persists the effective $(B,q)$ for reads and repair. The two object
strips illustrate different supported fixed layouts over equal byte extents,
not automatic-policy outcomes. The enlarged full-stripe shard inset illustrates
$s=B_1/d=2q_1$; packed frames and $q$ are distinct.
Dashed hint wiring still needs Entity/producer integration; smaller-$q$
writes require eligible new objects and both default-off gates enabled.
Upper reads verify whole objects before slicing. Other lookup paths,
Hole/visibility rules and lifecycle coordination are omitted.}

\vspace{14pt}
}]
\section{Private state over one shared snapshot}
Let $S$ be a pinned immutable snapshot and $\Delta_i$ the private upper of
workspace $i$. Its logical view is $W_i=\operatorname{Resolve}(S,\Delta_i)$.
The map stores explicit references to inherited ranges and to private data.
The lower groups names, common attributes and data locations into immutable
access structures. The figure follows one representative record into a
frame of $O_1$; it does not prescribe a fixed number of storage requests.

\section{A layout that stays with its object}
The coding block $B$ is requested when an object's finalized wire length is
known. RustFS admits the request, chooses a compatible verification quantum
$q$, and persists the effective layout. Existing objects are read and repaired
using their saved layout. The two illustrative strips show configurability,
not an adaptive optimization result. Integration and eligibility conditions
remain explicit in the caption.
\end{document}
